% Chapter Template

\chapter{Conclusiones} % Main chapter title

\label{Chapter5} % Change X to a consecutive number; for referencing this chapter elsewhere, use \ref{ChapterX}
Todos los capítulos deben comenzar con un breve párrafo introductorio que indique cuál es el contenido que se encontrará al leerlo.  La redacción sobre el contenido de la memoria debe hacerse en presente y todo lo referido al proyecto en pasado, siempre de modo impersonal.


%----------------------------------------------------------------------------------------

%----------------------------------------------------------------------------------------
%	SECTION 1
%----------------------------------------------------------------------------------------

\section{Conclusiones generales }

El trabajo aquí desarrollado logró cumplir con los objetivos planteados en la sección \ref{sec:objetivos}, ya que se agregó valor tecnológico al hogar y se integraron los conocimientos aprendidos en la Carrera de Especialidad en Sistemas Embebidos de la FIUBA.

Se cumplió con un 80\% de los requerimientos planteados en la planificación, resolviendo los relacionados al desarrollo del firmware, hardware y documentación, pero relegando a una segunda etapa la ejecución de la aplicación de usuario para smartphone y el manual de usuario. De esta manera, se manifestó el riesgo de incumplimiento de la planificación inicial, solventando el mismo con pequeñas alteraciones a fin de excluir esos requerimientos. Así se logró finalizar el sistema de iluminación antes de la fecha de entrega final a costo de incumplir el plan de mitigación, ya que no se logró extender los tiempos de desarrollo.

Fue útil para la ejecución de este trabajo disponer de los conocimientos técnicos adquiridos a lo largo de la carrera, así como aplicar lo planteado en la planificación. También fue de gran utilidad disponer de un director que acompañe con su sabiduría y disposición.

Sin embargo fue de poca utilidad la rigurosidad en la escritura de la memoria que obligó a disponer de mayor tiempo del planificado, quitando espacio al desarrollo técnico y tecnológico.


%----------------------------------------------------------------------------------------
%	SECTION 2
%----------------------------------------------------------------------------------------
\section{Próximos pasos}

Este proyecto deja abierto a un sin fin de posibilidades de mejoras y expansión. Algunos de los puntos a mejorar incluyen el desarrollo de la APP como interfaz responsiva y el manual de usuario. También el desarrollo de un PCB para el hardware, donde se incluyan todos los periféricos en una pequeña placa. Esto mejorará la estética del sistema, pero también la disposición de los componentes, disminuyendo el ruido sobre las señales, la disipación de potencia en forma de calor y mejorando el conexionado.

También se dejan abiertas las puertas para la implementación de un nuevo protocolo de comunicación entre el sistema y el usuario a través de un periférico WIFI, que mejore la conectividad, el alcance y la interacción con la persona que lo utilice. A su vez, se puede cambiar el microcontrolador de menor porte, a fin de reducir costos.